\chapter[ANALYSE ET CONCEPTION DU SYSTÈME]{ANALYSE ET CONCEPTION DU SYSTÈME}
    \section[Introduction partielle]{Introduction partielle}
    Ce chapitre vise à offrir une vue d‘ensemble sur l’analyse et une vue logique du
    système informatique qui sera développé. Dans le cycle de vie d’un système informatique,
    il existe deux étapes primordiales qui sont : l’analyse et la conception. L’analyse nous
    permet d’avoir une idée sur les besoins du travail. Il permet également d’avoir un aperçu
    du résultat. La conception, quant à elle, nous permet de décrire par un langage de
    modélisation de données (UML dans notre cas) la structure et le fonctionnement du
    système d’information.
    \section[Analyse des besoins et conception du système]{Analyse des besoins et conception du système}
        \subsection[La vue des besoins utilisateurs]{La vue des besoins utilisateurs}    
    \subsubsection[Identification des acteurs]{Identification des acteurs}
    Voici les acteurs identifiés dans le cadre de notre
    travail :
    \par
        \begin{itemize}
            \setlength{\itemsep}{0pt}
            \item [\ding{226}] \textbf{L’internaute} : est celui qui visite
            notre plateforme sans y être inscrit.
            \item [\ding{226}] \textbf{Le client} : est celui qui a créé un compte
            utilisateur dans le système.
            \item [\ding{226}] \textbf{L’administrateur} : il est chargé de
            monitorer le système afin d’éviter tout débordement des utilisateurs,
            et voir tous les mauvais fonctionnements du système et les régler
            à distance.
            \item [\ding{226}] \textbf{Le manager} : est celui qui utilise le système
            pour visualiser le tableau de bord et consulter le rapport.
        \end{itemize}
    \subsubsection[Identification des cas d’utilisation]{Identification des cas d’utilisation}
    Les cas d’utilisations ou fonctionnalités de notre
    système sont les suivants :
    \par
        \begin{itemize}
            \setlength{\itemsep}{0pt}
            \item [\ding{226}] Créer compte
            \item [\ding{226}] S’authentifier
            \item [\ding{226}] Consulter avis
            \item [\ding{226}] Afficher notification
            \item [\ding{226}] Scanner jeton
            \item [\ding{226}] Donner avis
            \item [\ding{226}] Établir rapport
            \item [\ding{226}] Vendre billet
            \item [\ding{226}] Gérer compte
            \item [\ding{226}] Ajouter évaluation
            \item [\ding{226}] Génère tableau de bord
            \item [\ding{226}] Consulter tableau de bord
            \item [\ding{226}] Consulter rapport
        \end{itemize}
    Chacun des cas figurant dans la liste présenter
    ci-haut, seras expliquer plus en détaille dans
    sa description textuelle.

%%%%%%%%%%%%%%%%%%%%%
%% Page en paysage %%
%%%%%%%%%%%%%%%%%%%%%
\uselandscape
    \subsubsection[Diagrammes de cas d’utilisation]{Diagrammes de cas d’utilisation}
    Après avoir eu à identifier les acteurs et leurs cas d’utilisation
    voici comment se présente le diagramme de cas d’utilisation :
        \begin{figure}[H]
            \centering
            \includegraphics[width=165mm]{images/dcu_systeme.png}
            \caption{Diagramme de cas d’utilisation}
            \label{fig:dcu}
        \end{figure}

%%%%%%%%%%%%%%%%%%%%%%
%% Page en portrait %%
%%%%%%%%%%%%%%%%%%%%%%        
\useportrait
    \subsubsection[Description textuelle]{Description textuelle}
        \paragraph[Créer compte]{Créer compte}
            \begin{longtable}{p{4cm} p{9cm}}
                \caption{Description textuelle du cas d’utilisation Créer compte}
                \label{table:usecaseCreeCompte}
                \\\hline\hline
                    \textbf{Cas d’utilisation} & \textbf{Créer compte}
                \\\hline\hline
                        \textbf{Acteur principal} & L’\textsc{internaute}
                    \\
                        \textbf{Acteur secondaire} & -
                    \\
                        \textbf{Objectifs} & L’\textsc{internaute} veut être 
                        enregistré dans le système.
                    \\
                        \textbf{Préconditions} & -
                    \\
                        \textbf{Postcondition} & Un compte a était créé.
                    \\
                        \textbf{Scénario nominal} &
                            \begin{enumerate}[leftmargin=*]
                                \item L’\textsc{internaute} clique sur créer compte et le
                                système affiche un formulaire.
                                \item L’\textsc{internaute} saisi ses identifiants à savoir
                                (nom, post-nom, date et année de
                                naissance, âge et genre) et coordonnées (téléphone, adresse physique
                                et adresse électronique).
                                \item L’\textsc{internaute} confirme en cliquant sur créer.
                                \item Le système enregistre les nouvelles informations
                                dans la base de données et le
                                système répond avec un message de succès.
                            \end{enumerate}
                    \\
                        \textbf{Scénario alternatif} &
                            \begin{enumerate}[leftmargin=*]
                                \item À partir de l’étape numéro 3 du scenario nominal
                                le système vérifie et ne valide pas le formulaire
                                \item Retour à l’étape numéro 1 du scénario principal
                            \end{enumerate}
                \\\bottomrule
            \end{longtable}
            %\pagebreak

        \paragraph[S’authentifier]{S’authentifier}
            \begin{longtable}{p{4cm} p{9cm}}
                \caption{Description textuelle du cas d’utilisation S’authentifier}
                \label{table:usecaseSauth}
                \\\hline\hline
                    \textbf{Cas d’utilisation} & \textbf{S’authentifier}
                \\\hline\hline
                        \textbf{Acteur principal} & L’\textsc{internaute}, le \textsc{client},
                        le \textsc{guichetier}, le \textsc{manager} et l’\textsc{administrateur}
                    \\
                        \textbf{Acteur secondaire} & -
                    \\
                        \textbf{Objectifs} & L’utilisateur veut accéder au système.
                    \\
                        \textbf{Préconditions} & -
                    \\
                    \textbf{Scénario nominal} &
                        \begin{enumerate}[leftmargin=*]
                            \item L’utilisateur saisis les informations de connexion à savoir (le mot de
                            passe et l’identifient).
                            \item Il confirme en cliquant sur le bouton se connecter.
                            \item Le système vérifie les informations saisies et crée la session.
                        \end{enumerate}
                    \\
                    \textbf{Scénario alternatif} &
                    \begin{enumerate}[leftmargin=*]
                            \item mm
                            \item kk
                        \end{enumerate}
                    \\
                    \textbf{Postcondition}
                \\\bottomrule
            \end{longtable}

        \paragraph[Consulter avis]{Consulter avis}
            \begin{longtable}{p{4cm} p{9cm}}
                \caption{Description textuelle du cas d’utilisation Consulter avis}
                \label{table:usecaseConsulterAvis}
                \\\hline\hline
                    \textbf{Cas d’utilisation} & \textbf{Consulter avis}
                \\\hline\hline
                        \textbf{Acteur principal} & aa
                    \\
                        \textbf{Acteur secondaire} & bb
                    \\
                        \textbf{Objectifs} & dd
                    \\
                        \textbf{Préconditions} & dd
                    \\
                    \textbf{Scénario nominal} &
                        \begin{enumerate}[leftmargin=*]
                            \item dd
                            \item ddd
                        \end{enumerate}
                    \\
                    \textbf{Scénario alternatif} &
                        \begin{enumerate}[leftmargin=*]
                            \item mm
                            \item kk
                        \end{enumerate}
                    \\
                    \textbf{Postcondition}
                \\\bottomrule
            \end{longtable}

        \paragraph[Afficher notification]{Afficher notification}
        \begin{longtable}{p{4cm} p{9cm}}
            \caption{Description textuelle du cas d’utilisation Afficher notification}
            \label{table:usecaseAfficherNotification}
            \\\hline\hline
                \textbf{Cas d’utilisation} & \textbf{Afficher notification}
            \\\hline\hline
                    \textbf{Acteur principal} & aa
                \\
                    \textbf{Acteur secondaire} & bb
                \\
                    \textbf{Objectifs} & dd
                \\
                    \textbf{Préconditions} & dd
                \\
                \textbf{Scénario nominal} &
                    \begin{enumerate}[leftmargin=*]
                        \item dd
                        \item ddd
                    \end{enumerate}
                \\
                \textbf{Scénario alternatif} &
                    \begin{enumerate}[leftmargin=*]
                        \item mm
                        \item kk
                    \end{enumerate}
                \\
                \textbf{Postcondition}
            \\\bottomrule
        \end{longtable}

        \paragraph[Scanner jeton]{Scanner jeton}
        \begin{longtable}{p{4cm} p{9cm}}
            \caption{Description textuelle du cas d’utilisation Scanner jeton}
            \label{table:usecaseScannerJ}
            \\\hline\hline
                \textbf{Cas d’utilisation} & \textbf{Scanner jeton}
            \\\hline\hline
                    \textbf{Acteur principal} & aa
                \\
                    \textbf{Acteur secondaire} & bb
                \\
                    \textbf{Objectifs} & dd
                \\
                    \textbf{Préconditions} & dd
                \\
                \textbf{Scénario nominal} &
                    \begin{enumerate}[leftmargin=*]
                        \item dd
                        \item ddd
                    \end{enumerate}
                \\
                \textbf{Scénario alternatif} &
                    \begin{enumerate}[leftmargin=*]
                        \item mm
                        \item kk
                    \end{enumerate}
                \\
                \textbf{Postcondition}
            \\\bottomrule
        \end{longtable}

        \paragraph[Donner avis]{Donner avis}
        \begin{longtable}{p{4cm} p{9cm}}
            \caption{Description textuelle du cas d’utilisation Donner avis}
            \label{table:usecaseDonnerAvis}
            \\\hline\hline
                \textbf{Cas d’utilisation} & \textbf{Donner avis}
            \\\hline\hline
                    \textbf{Acteur principal} & aa
                \\
                    \textbf{Acteur secondaire} & bb
                \\
                    \textbf{Objectifs} & dd
                \\
                    \textbf{Préconditions} & dd
                \\
                \textbf{Scénario nominal} &
                    \begin{enumerate}[leftmargin=*]
                        \item dd
                        \item ddd
                    \end{enumerate}
                \\
                \textbf{Scénario alternatif} &
                    \begin{enumerate}[leftmargin=*]
                        \item mm
                        \item kk
                    \end{enumerate}
                \\
                \textbf{Postcondition}
            \\\bottomrule
        \end{longtable}

        \paragraph[Établir rapport]{Établir rapport}
        \begin{longtable}{p{4cm} p{9cm}}
            \caption{Description textuelle du cas d’utilisation Établir rapport}
            \label{table:usecaseEtablirCompte}
            \\\hline\hline
                \textbf{Cas d’utilisation} & \textbf{Établir rapport}
            \\\hline\hline
                    \textbf{Acteur principal} & aa
                \\
                    \textbf{Acteur secondaire} & bb
                \\
                    \textbf{Objectifs} & dd
                \\
                    \textbf{Préconditions} & dd
                \\
                \textbf{Scénario nominal} &
                    \begin{enumerate}[leftmargin=*]
                        \item dd
                        \item ddd
                    \end{enumerate}
                \\
                \textbf{Scénario alternatif} &
                    \begin{enumerate}[leftmargin=*]
                        \item mm
                        \item kk
                    \end{enumerate}
                \\
                \textbf{Postcondition}
            \\\bottomrule
        \end{longtable}

        \paragraph[Vendre billet]{Vendre billet}
        \begin{longtable}{p{4cm} p{9cm}}
            \caption{Description textuelle du cas d’utilisation Vendre billet}
            \label{table:usecaseVendreBill}
            \\\hline\hline
                \textbf{Cas d’utilisation} & \textbf{Vendre billet}
            \\\hline\hline
                    \textbf{Acteur principal} & aa
                \\
                    \textbf{Acteur secondaire} & bb
                \\
                    \textbf{Objectifs} & dd
                \\
                    \textbf{Préconditions} & dd
                \\
                \textbf{Scénario nominal} &
                    \begin{enumerate}[leftmargin=*]
                        \item dd
                        \item ddd
                    \end{enumerate}
                \\
                \textbf{Scénario alternatif} &
                    \begin{enumerate}[leftmargin=*]
                        \item mm
                        \item kk
                    \end{enumerate}
                \\
                \textbf{Postcondition}
            \\\bottomrule
        \end{longtable}

        \paragraph[Gérer compte]{Gérer compte}
        \begin{longtable}{p{4cm} p{9cm}}
            \caption{Description textuelle du cas d’utilisation Gérer compte}
            \label{table:usecaseGereCompte}
            \\\hline\hline
                \textbf{Cas d’utilisation} & \textbf{Gérer compte}
            \\\hline\hline
                    \textbf{Acteur principal} & aa
                \\
                    \textbf{Acteur secondaire} & bb
                \\
                    \textbf{Objectifs} & dd
                \\
                    \textbf{Préconditions} & dd
                \\
                \textbf{Scénario nominal} &
                    \begin{enumerate}[leftmargin=*]
                        \item dd
                        \item ddd
                    \end{enumerate}
                \\
                \textbf{Scénario alternatif} &
                    \begin{enumerate}[leftmargin=*]
                        \item mm
                        \item kk
                    \end{enumerate}
                \\
                \textbf{Postcondition}
            \\\bottomrule
        \end{longtable}

        \paragraph[Ajouter évaluation]{Ajouter évaluation}
        \begin{longtable}{p{4cm} p{9cm}}
            \caption{Description textuelle du cas d’utilisation Ajouter évaluation}
            \label{table:usecaseAjouterEval}
            \\\hline\hline
                \textbf{Cas d’utilisation} & \textbf{Ajouter évaluation}
            \\\hline\hline
                    \textbf{Acteur principal} & aa
                \\
                    \textbf{Acteur secondaire} & bb
                \\
                    \textbf{Objectifs} & dd
                \\
                    \textbf{Préconditions} & dd
                \\
                \textbf{Scénario nominal} &
                    \begin{enumerate}[leftmargin=*]
                        \item dd
                        \item ddd
                    \end{enumerate}
                \\
                \textbf{Scénario alternatif} &
                    \begin{enumerate}[leftmargin=*]
                        \item mm
                        \item kk
                    \end{enumerate}
                \\
                \textbf{Postcondition}
            \\\bottomrule
        \end{longtable}

        \paragraph[Générer tableau de bord]{Générer tableau de bord}
        \begin{longtable}{p{4cm} p{9cm}}
            \caption{Description textuelle du cas d’utilisation Générer tableau de bord}
            \label{table:usecaseDashboard}
            \\\hline\hline
                \textbf{Cas d’utilisation} & \textbf{Générer tableau de bord}
            \\\hline\hline
                    \textbf{Acteur principal} & aa
                \\
                    \textbf{Acteur secondaire} & bb
                \\
                    \textbf{Objectifs} & dd
                \\
                    \textbf{Préconditions} & dd
                \\
                \textbf{Scénario nominal} &
                    \begin{enumerate}[leftmargin=*]
                        \item dd
                        \item ddd
                    \end{enumerate}
                \\
                \textbf{Scénario alternatif} &
                    \begin{enumerate}[leftmargin=*]
                        \item mm
                        \item kk
                    \end{enumerate}
                \\
                \textbf{Postcondition}
            \\\bottomrule
        \end{longtable}

        \paragraph[Consulter tableau de bord]{Consulter tableau de bord}
        \begin{longtable}{p{4cm} p{9cm}}
            \caption{Description textuelle du cas d’utilisation Consulter tableau de bord}
            \label{table:usecaseConsDash}
            \\\hline\hline
                \textbf{Cas d’utilisation} & \textbf{Consulter tableau de bord}
            \\\hline\hline
                    \textbf{Acteur principal} & aa
                \\
                    \textbf{Acteur secondaire} & bb
                \\
                    \textbf{Objectifs} & dd
                \\
                    \textbf{Préconditions} & dd
                \\
                \textbf{Scénario nominal} &
                    \begin{enumerate}[leftmargin=*]
                        \item dd
                        \item ddd
                    \end{enumerate}
                \\
                \textbf{Scénario alternatif} &
                    \begin{enumerate}[leftmargin=*]
                        \item mm
                        \item kk
                    \end{enumerate}
                \\
                \textbf{Postcondition}
            \\\bottomrule
        \end{longtable}

        \paragraph[Consulter rapport]{Consulter rapport}
        \begin{longtable}{p{4cm} p{9cm}}
            \caption{Description textuelle du cas d’utilisation Consulter rapport}
            \label{table:usecaseConsulterRapport}
            \\\hline\hline
                \textbf{Cas d’utilisation} & \textbf{Consulter rapport}
            \\\hline\hline
                    \textbf{Acteur principal} & aa
                \\
                    \textbf{Acteur secondaire} & bb
                \\
                    \textbf{Objectifs} & dd
                \\
                    \textbf{Préconditions} & dd
                \\
                \textbf{Scénario nominal} &
                    \begin{enumerate}[leftmargin=*]
                        \item dd
                        \item ddd
                    \end{enumerate}
                \\
                \textbf{Scénario alternatif} &
                    \begin{enumerate}[leftmargin=*]
                        \item mm
                        \item kk
                    \end{enumerate}
                \\
                \textbf{Postcondition}
            \\\bottomrule
        \end{longtable}
        \input{chapitres/03-Chap-3.2/2-vue-du-processus.tex}
        \input{chapitres/03-Chap-3.2/3-vue-logique.tex}
    \section[Conclusion partielle]{Conclusion partielle}